\chapter[致谢]{致\quad 谢}\chaptermark{致\quad 谢}% syntax: \chapter[目录]{标题}\chaptermark{页眉}
% \thispagestyle{noheaderstyle}% 如果需要移除当前页的页眉
%\pagestyle{noheaderstyle}% 如果需要移除整章的页眉

\reviewORprint{
}{
    时光飞逝，思绪万千，两年前我完成本科毕业设计，两年后我已经是一位即将进入博士生涯的硕士二年级学生了。两年的时间里，我成为了一名极度实用主义者。我学习鼓捣实验室的各种设备，显卡集群，网络节点。

    首先感谢陈东岳老师和邓诗卓老师的辛勤指导，两位老师培养我的学术品味，教会我学术思维。每一次组会都是思想火花的碰撞，都会迸发新的学术idea。硕士毕业之后我还继续留在课题组里攻读博士学位，会继续接受老师们的指导。

    云山苍苍，江水泱泱，先生之风，山高水长。
    
    其次感谢朱越师兄，教会了我很多技能。朱越师兄幽默风趣，学术知识渊博，时时刻刻都在熏陶我，在此祝朱越师兄博士毕业能年薪百万！

    接下来感谢我的球友们，我们在挥拍时放松自我，在运动场上迸发青春的活力，他们是：张宸宇，余伟杰，杜鑫宇，徐宁，蒋家首都，张哲睿，林奔，武昊阳，江楚玥，季鹏，张曦文（上述排名不分先后），还有很多伙伴们，在此我就不一一列举了，祝你们学习进步，工作顺利！

    感谢十舍A区319宿舍的三位舍友们，李哲瀚，魏俊杰，王敬博，我们一起生活，一起面对挫折，一起成长，祝你们都能找到理想的工作。

    感谢同门师兄师姐师妹师弟的陪伴和照顾，祝你们学习进步，工作顺利！

    感谢父母的辛勤栽培，祝你们身体健康，工作顺利。

    感谢七年前选择复读的自己，不然我就得去学医了。今天人工智能发展如火如荼，我想投身于其浪潮之中，希望有朝一日可以见证AGI的诞生。

    衷心祝愿生命中的遇到的每一个人，身体健康，工作顺利，快乐多多，烦恼少少。
}

\chapter{攻读硕士学位期间取得的学术成果}

\reviewORprint{
    \section*{学术论文:}
    \begin{enumerate}[leftmargin=*]
        \item Deng S, Yang Q, Yang Z, et al. Tri-modality Collaborative Learning for Person Re-identification[C]//Australasian Database Conference. Singapore: Springer Nature Singapore, 2024: 322--333.（EI）
        \item Yang Q, Deng S, Chen D, et al. FRWKV: Frequency-Domain Linear Attention for Long-Term Time Series Forecasting[J]. arXiv preprint arXiv:2512.07539, 2025.
    \end{enumerate}
}{
    \section*{学术论文:}
    \begin{enumerate}[leftmargin=*]
        \item Deng S, Yang Q, Yang Z, et al. Tri-modality Collaborative Learning for Person Re-identification[C]//Australasian Database Conference. Singapore: Springer Nature Singapore, 2024: 322--333.（EI）
        \item Yang Q, Deng S, Chen D, et al. FRWKV: Frequency-Domain Linear Attention for Long-Term Time Series Forecasting[J]. arXiv preprint arXiv:2512.07539, 2025.
    \end{enumerate}
}

\cleardoublepage[plain]% 让文档总是结束于偶数页，可根据需要设定页眉页脚样式，如 [noheaderstyle]
